\documentclass[12pt,a4paper]{article}
\usepackage[utf8]{inputenc}
\usepackage[turkish]{babel}
\usepackage{amsmath}
\usepackage{graphicx}
\usepackage{hyperref}
\usepackage{listings}
\usepackage{xcolor}
\usepackage{geometry}
\usepackage{enumitem}
\geometry{a4paper, margin=1in}

\definecolor{codegreen}{rgb}{0,0.6,0}
\definecolor{codegray}{rgb}{0.5,0.5,0.5}
\definecolor{codepurple}{rgb}{0.58,0,0.82}
\definecolor{backcolour}{rgb}{0.95,0.95,0.92}

\lstdefinestyle{mystyle}{
    backgroundcolor=\color{backcolour},   
    commentstyle=\color{codegreen},
    keywordstyle=\color{magenta},
    numberstyle=\tiny\color{codegray},
    stringstyle=\color{codepurple},
    basicstyle=\ttfamily\footnotesize,
    breakatwhitespace=false, breaklines=true, captionpos=b,
    keepspaces=true, numbers=left, numbersep=5pt,
    showspaces=false, showstringspaces=false, showtabs=false, tabsize=2
}
\lstset{style=mystyle}

\begin{document}
\shorthandoff{=}

\begin{titlepage}
    \centering
    \vspace*{1cm}
    {\Large\textbf{ANKARA UNIVERSITESI}}\\[0.3cm]
    {\large Muhendislik Fakultesi}\\[0.3cm]
    {\large Bilgisayar Muhendisligi Bolumu}\\[1cm]
    \begin{figure}[h]
    \centering
    \vspace{3pt}
    \includegraphics[width=5.0cm]{logo.png}
    \end{figure}
    \begin{tabular}{ll}
        \textbf{Ogrenci:} & Ahmet Akgun \\
        \textbf{Ogrenci No:} & 22290602 \\
        \textbf{Danisman:} & Enver Bagci \\
    \end{tabular}\\[1.5cm]
    {\Large \textbf{BLM4522 Ag Tabanli Paralel Dagitim Sistemleri}}\\[0.5cm]
    {\large \textbf{Proje Adi:} Veri Temizleme ve ETL Surecleri Tasarimi}
    \vfill
    \begin{itemize}[label=\textbullet]
        \item \textbf{GitHub Depo Linki:} \url{https://github.com/aakgunnn}
        \item \textbf{Haftalik Video Takip Baglantisi:} \url{https://www.youtube.com/@ahmetakgun359}
    \end{itemize}
\end{titlepage}

\tableofcontents
\newpage

\section{Giris}
Bu proje kapsaminda, buyuk veri kumelerinin temizlenmesi ve islenmesi icin bir ETL (Extract, Transform, Load) pipeline'i PostgreSQL uzerinde tamamen SQL tabanli olarak tasarlanmis ve uygulanmistir. Pipeline uc SQL script dosyasindan olusmaktadir:

\begin{enumerate}
    \item \textbf{stage1\_extract.sql:} Ham CSV dosyalarindan veriyi staging tablolarina yukler.
    \item \textbf{stage2\_transform.sql:} Veriyi temizler, dogrular, standartlastirir ve hedef tablolara yukler.
    \item \textbf{stage3\_quality\_report.sql:} Veri kalitesi raporlari olusturur.
\end{enumerate}

\section{Ham Veri Kaynaklari}
Projede iki adet ham CSV dosyasi kullanilmistir. Bu dosyalar kasitli olarak cesitli veri kalitesi sorunlari icerecek sekilde tasarlanmistir:

\subsection{Musteri Verisi (raw\_data\_customers.csv)}
13 satirlik bu dosyada: eksik isim alanlari, tutarsiz email formatlari, farkli telefon ve tarih formatlari, tutarsiz sehir adlari, negatif/NULL tutarlar, duplicate kayitlar ve fazla bosluklar bulunmaktadir.

\subsection{Siparis Verisi (raw\_data\_orders.csv)}
12 satirlik bu dosyada: negatif/sifir miktar degerleri, sayisal olmayan miktar (\texttt{abc}), eksik fiyat bilgileri, sistemde olmayan musteri ID'leri, tutarsiz siparis durumlari ve tekrarlanan siparisler bulunmaktadir.

\section{Extract (Veri Cikarma) -- stage1\_extract.sql}
Pipeline'in ilk asamasinda, ham veriler tum sutunlari \texttt{VARCHAR} tipinde tutan \textbf{staging tablolarina} yuklenmistir. Boylece henuz dogrulanmamis verilerin tip hatalari olusturmasi engellenmistir.

\begin{lstlisting}[language=SQL, caption=Staging Tablolari ve Veri Yukleme]
CREATE TABLE staging_customers (
    id VARCHAR(10),
    full_name VARCHAR(200),
    email VARCHAR(200),
    phone VARCHAR(50),
    city VARCHAR(100),
    registration_date VARCHAR(50),
    total_spent VARCHAR(50)
);

\copy staging_customers FROM 'raw_data_customers.csv' WITH (FORMAT csv, HEADER true);
\end{lstlisting}

\section{Transform (Veri Donusturme) -- stage2\_transform.sql}
Donusturme asamasi, pipeline'in en kapsamli bolumdur. Tum temizleme ve standartlastirma islemleri tek bir SQL \texttt{INSERT INTO ... SELECT} ifadesi icinde gerceklestirilmistir.

\subsection{Email Dogrulama}
Regex tabanli email dogrulamasi ile gecersiz formattaki adresler (\texttt{elif\_celik@company}, \texttt{can.polat@}) tespit edilerek ilgili kayitlar atlanmistir:
\begin{lstlisting}[language=SQL, caption=Regex ile Email Dogrulama]
WHERE TRIM(email) ~ '^[A-Za-z0-9._%+-]+@[A-Za-z0-9.-]+\.[A-Za-z]{2,}$'
\end{lstlisting}

\subsection{Telefon Standardizasyonu}
Farkli formatlardaki telefon numaralari (\texttt{05551234568} vs \texttt{+90 555 123 4567}) SQL \texttt{CASE} ifadeleri ve \texttt{REGEXP\_REPLACE} fonksiyonu ile tek bir standart formata donusturulmustur:
\begin{lstlisting}[language=SQL, caption=Telefon Format Donusumu]
CASE
  WHEN LENGTH(REGEXP_REPLACE(phone, '\D', '', 'g')) = 11 
    THEN '+90 ' || SUBSTRING(REGEXP_REPLACE(phone, '\D', '', 'g'), 2, 3)
         || ' ' || SUBSTRING(..., 5, 3)
         || ' ' || SUBSTRING(..., 8, 4)
  ...
END AS phone
\end{lstlisting}

\subsection{Tarih ve Sehir Standardizasyonu}
\texttt{YYYY-MM-DD}, \texttt{DD/MM/YYYY} ve \texttt{YYYY/MM/DD} gibi farkli tarih formatlari regex pattern eslesimiyle otomatik algilanip \texttt{TO\_DATE} fonksiyonu ile ISO 8601 standardina cevrilmistir. Sehir adlari \texttt{INITCAP()} fonksiyonu ile standartlastirilmistir (\texttt{ankara} $\rightarrow$ \texttt{Ankara}).

\subsection{Negatif ve NULL Deger Kontrolu}
\begin{lstlisting}[language=SQL, caption=Negatif/NULL Tutar Duzeltme]
GREATEST(
    COALESCE(
        CASE 
            WHEN UPPER(TRIM(total_spent)) = 'NULL' OR TRIM(total_spent) = '' THEN 0
            ELSE CAST(total_spent AS NUMERIC(12,2))
        END, 0
    ), 0
) AS total_spent
\end{lstlisting}

\subsection{Duplicate Kontrolu}
\texttt{DISTINCT ON} ifadesi ile ayni email adresine sahip tekrarlanan kayitlardan sadece ilki alinmistir:
\begin{lstlisting}[language=SQL, caption=Duplicate Eleme]
SELECT DISTINCT ON (LOWER(TRIM(email))) *
FROM staging_customers
ORDER BY LOWER(TRIM(email)), id::INTEGER
\end{lstlisting}

\section{Load (Veri Yukleme)}
Temizlenmis veriler, referans butunlugu (REFERENCES) ve is kurallari (CHECK) kisitlarina sahip hedef tablolara yuklenistir:
\begin{lstlisting}[language=SQL, caption=Hedef Tablo Semalari]
CREATE TABLE clean_customers (
    customer_id SERIAL PRIMARY KEY,
    full_name VARCHAR(100) NOT NULL,
    email VARCHAR(100) UNIQUE NOT NULL,
    phone VARCHAR(20),
    city VARCHAR(50),
    registration_date DATE,
    total_spent NUMERIC(12,2) DEFAULT 0.00
);

CREATE TABLE clean_orders (
    order_id INTEGER PRIMARY KEY,
    customer_id INTEGER REFERENCES clean_customers(customer_id),
    quantity INTEGER CHECK (quantity > 0),
    ...
);
\end{lstlisting}

\section{Veri Kalitesi Raporu -- stage3\_quality\_report.sql}
Son asamada, surece dair kapsamli bir denetim raporu SQL sorgulari ile olusturulmustur. Rapor su bilgileri icermektedir:

\begin{itemize}
    \item Ham vs. temiz kayit sayisi karsilastirmasi
    \item Atlanan musterilerin detayli listesi (neden + hangi satir)
    \item Atlanan siparislerin detayli listesi
    \item Donusturme detaylari (negatif tutar, NULL deger duzeltmeleri)
    \item Temiz veri dogrulama (tum temiz kayitlarin listesi)
    \item Sehir bazli musteri dagilimi
    \item Durum bazli siparis ve ciro istatistikleri
\end{itemize}

\begin{lstlisting}[caption=Rapor Ciktisi (Ozet)]
        metrik        | deger
----------------------+-------
 Ham Musteri Sayisi   | 13
 Temiz Musteri Sayisi | 9
 Ham Siparis Sayisi   | 12
 Temiz Siparis Sayisi | 6

          metrik          |    deger
--------------------------+--------------
 TOPLAM CIRO (Tamamlandi) | 18,000.00 TL
\end{lstlisting}

\section{Sonuc}
Bu projede, gercek dunya senaryolarina uygun bir ETL pipeline'i tamamen SQL tabanli olarak tasarlanmis ve PostgreSQL uzerinde basariyla uygulanmistir. Regex (\texttt{\~}), \texttt{CASE}, \texttt{COALESCE}, \texttt{GREATEST}, \texttt{INITCAP}, \texttt{DISTINCT ON}, \texttt{TO\_DATE} ve \texttt{REGEXP\_REPLACE} gibi guclu SQL fonksiyonlari kullanilarak veri temizleme, dogrulama, standartlastirma ve raporlama islemleri gerceklestirilmistir.

\end{document}
